\documentclass[a4paper,12pt]{article}
\usepackage[spanish]{babel}
\usepackage[dvipdfm,dvipsnames,usenames]{color}
\usepackage{amsmath,amsfonts}
\usepackage{eurosym}
\usepackage[utf8]{inputenc}


\begin{document}

\leftline{\bf Introducci\'{o}n} 

Aqu\'i empieza la introducci\'on.\\
Estos son ejemplos de TIPOS de letra \\
{\bf bold}\textbf{face}\\
{\rm ro}\textrm{man}\\
{\sl slan}\textsl{ted}\\
{\sf sans} \textsf{serif}\\
{\sc small} \textsc{caps}\\
{\tt type}\texttt{writer}\\
\chapter{Estado del arte}
Aqu\'i empieza el cap\'itulo sobre estado del arte\\
OTROS EJEMPLOS \\
{\tiny Que} {\scriptsize cant} {\footnotesize id} {\small ad} {\normalsize de} {\large tam} {\Large anos} {\LARGE de} {\huge le} {\Huge tra}

\begin{center}
\huge{examen de teoría} \\
\bigskip
Primer curso \\
Enero del 2019
\end{center}

\centerline{\bf línea centrada negrita}

En el ambiente \textit{verbatim}, \textbf{LATEX} procesa el texto exactamente como está escrito utilizando fuente \textit{typewriter}. Util para secciones de \textit{código en C o Fortran}.

\begin{verbatim}
c bucle en %i
do i=1, n
a(i,i+1) = i
end do
\end{verbatim}




\end{document}
